\documentclass[10pt,a4paper]{amsart}
\usepackage[margin=19mm]{geometry}
\usepackage{amsmath,amssymb,mathtools}
\usepackage[scr=rsfs,cal=euler]{mathalfa}
\usepackage{xfrac}
\usepackage[unicode,hidelinks,bookmarks=false]{hyperref}
\newcommand{\ie}{i.e.\ }
\newcommand{\cf}{cf.\ }
\newcommand{\resp}{respectively\ }
\newcommand{\Subsection}{\subsection}
\providecommand{\nobreakdash}{\nobreak\hbox{-}}
\setlength{\emergencystretch}{2em}
\multlinegap0pt
\usepackage{amssymb}
\usepackage[shortlabels]{enumitem}
\usepackage{float}
\usepackage{mathtools}
\newtheorem{lem}{Lemma}
\newcommand{\Z}{\mathbb{Z}}
\title{Decent actions of groups on restricted products}
\author{Chris Karpinski}
\date{Source: arXiv:2604.22635v1 (CC BY 4.0)}
\begin{document}
\maketitle

\noindent\emph{Source and licence.} Decent actions of groups on restricted products. Version arXiv:2604.22635v1 (CC BY 4.0), distributed by arXiv under the Creative Commons Attribution 4.0 International licence, http://creativecommons.org/licenses/by/4.0/, as recorded in the arXiv licence metadata for this exact version and shown on the abstract page https://arxiv.org/abs/2604.22635v1. Full licence notice: \texttt{licenses/arxiv-2604.22635-CC-BY-4.0.txt}.\par\smallskip

\noindent\textit{Editorial scope.} Counted unit: the article's lem pt (source0.tex lines 222--227). The article's complete original proof of it is delivered with it (source0.tex lines 229--249, one \texttt{proof} environment of 21 lines). No mathematics is written, repaired or re-derived: the statement and the proof are the article's own lines, in the article's own notation, and no retained line is substituted or re-flowed.  No further source-local context is needed: every cross-reference of every retained line resolves inside the two delivered spans.  Nothing was omitted from any retained span, so the package carries no omission record.  No retained line contains a citation, so the wrapper carries no bibliography and no bibliography entry.  No retained line contains a figure environment, an image-inclusion command or drawing code, so no image is delivered and the package declares no asset dependency.  Editorial formatting only, in addition: an AMS \texttt{amsart} preamble that restates the article's own declarations and macros used by the retained lines, namely the environments \texttt{lem} on separate counters, together with the standard packages those lines need.  The retained environments print in this document as Lemma 1 in order of appearance, whereas the article prints the same items with the numbering of its own full document.  The complete native source of the article as distributed in the arXiv e-print for this exact version is bundled unchanged as \texttt{sources/199114/source0.tex}.
\begin{lem}
    \label{Inf orb fix pt}

    Let $t \in G^{\oplus}$ and let $O$ be an infinite orbit of $t$ on $P$. Then for any fixed points $y,z$ of $t$ in $X$, we have $y_p = z_p$ for each $p \in O$. 
    
\end{lem}

\begin{proof}
    Let $z$ be a fixed point of $t = gh$ in $X$ for $g \in \prod_P G_0$ and $h \in H$. We will show that the coordinates of $z$ in $O$ are uniquely determined. 

    Fixing $p \in O$, we can enumerate the elements of $O$ as $O = \{t^n p : n \in \Z\}$. Since $O$ is infinite, the map $\Z \ni n \mapsto t^n p$ is a bijection, we can identify $O$ with $\Z$ and write the $t$-action on $O$ as $t^n i = i + n$ for any $i,n \in \Z$. We will also write $z_{t^n p} (= z_{h^n p})$ as $z_n$ and $g_{t^n p} (= g_{h^n p})$ as $g_n$.

    Since $t(z) = z$, we have for each $n \in \Z$ that $g_n z_{n-1} = z_n$. Denote $S = \mathrm{Stab}_{G_0}(x_0)$. Since $t \in G^{\oplus}$, there exists $N \leq M$ in $\Z$ such that $g_n \in S$ for all $n \leq N$ and for all $n \geq M$. Up to re-indexing $O$, we can assume that $N = 0$. 

    We claim that $z_n = x_0$ for all $n \leq 0$. Indeed, since $z \in X$, there exists $K \leq 0$ such that for all $n \leq K$, $z_n = x_0$. By induction, for each $K \leq n \leq 0$, we have that 
    
    $$z_n = g_nz_{n-1} = g_n g_{n-1} z_{n-2} = \cdots = g_n g_{n-1} \cdots g_K z_K$$
    
    Since $g_i \in S$ for all $i \leq 0$, the latter expression is in $S \{z_K\} = S \{x_0\} = \{x_0\}$. Therefore, $z_n = x_0$ for all $K \leq n \leq 0$, and hence for all $n \leq 0$ (since $K$ was chosen such that $z_n = x_0$ for all $n \leq K$). Note that a symmetric argument yields that $z_n = x_0$ for all $n \geq M$, but we will not need this.

    Now for each $n > 0$, by induction we have: 

    $$z_n = g_n z_{n-1} = g_n g_{n-1} z_{n-2} = \cdots = g_n g_{n-1} \cdots g_1 z_0 = g_n g_{n-1} \cdots g_1 x_0$$

    Thus, the coordinates of $z$ on $O$ depend only on $g$ and $x_0$, hence are uniquely determined. 

    
\end{proof}

\end{document}
